\ifdefined\gkver\else\def\gkver{student}\fi
\documentclass[\gkver]{gaokaozhenti}
\begin{document}

\chapter{2007年海南理综}

\begin{enumerate}
\item
%% number: 1
%% paperName: 2007年海南理综第 1 题 3 分
%% typeId: 1
%% type: 单选题
%% chapter: 运动和力
%% point: 牛顿第一定律
%% method: 无
%% score: 3
%% degree: 650
%% duplicateId: 0
%% body:
16 世纪纪末，伽利略用实验和推理，推翻了已在欧洲流行了近两千年的亚里士多德关于力和运动的理论，开启了物理学发展的新纪元。在以下说法中，与亚里士多德观点相反的是\xzanswer{D}

\fourchoices[answer=D]
{四匹马拉拉车比两匹马拉的车跑得快：这说明，物体受的力越大，速度就越大}
{一个运动的物体，如果不再受力了，它总会逐渐停下来，这说明，静止状态才是物体长时间不受力时的“自然状态”}
{两物体从同一高度自由下落，较重的物体下落较快}
{一个物体维持匀速直线运动，不需要受力}

%% answer: D
%% memo:
\memoanswer{
亚里士多德认为，力是维持物体运动的原因，物体受力运动，不受力停止。伽利略通过理想斜面实验，说明了物体在不受力作用时总保持原来的运动状态，即力是改变物体运动状态的原因，故 D 正确。
}

\item
%% number: 2
%% paperName: 2007年海南理综第 2 题 3 分
%% typeId: 1
%% type: 单选题
%% chapter: 力物体的平衡
%% point: 力
%% method: 无
%% score: 3
%% degree: 850
%% duplicateId: 0
%% body:
如图，P 是位于水平的粗糙桌面上的物块。用跨过定滑轮的轻绳将 P 与小盘相连，小盘内有硅码，小盘与硅码的总质量为 $m$。在 P 运动的过程中，若不计空气阻力，则关于 P 在水平面方向受到的作用力与相应的施力物体，下列说法正确的是\xzanswer{B}

\onepicture[w=5.0cm]{\includesvg{figs/02}}

\fourchoices[answer=B]
{拉力和摩擦力，施力物体是地球和桌面}
{拉力和摩擦力，施力物体是绳和桌面}
{重力 $mg$ 和摩擦力，施力物体是地球和桌面}
{重力 $mg$ 和摩擦力，施力物体是绳和桌面}

%% answer: B
%% memo:
\memoanswer{
P 在水平方向受到绳子对它的拉力和桌子对它的摩擦力的作用，故 B 正确。
}

\item
%% number: 3
%% paperName: 2007年海南理综第 3 题 3 分
%% typeId: 1
%% type: 单选题
%% chapter: 电磁感应
%% point: 自感与互感，涡流
%% method: 无
%% score: 3
%% degree: 650
%% duplicateId: 0
%% body:
在如图所示的电路中，a、b 为两个完全相同的灯泡，$L$ 为自感线圈，$E$ 为电源，S 为开关。关于两灯泡点亮和熄灭的先后次序，下列说法正确的是\xzanswer{C}

\onepicture[w=3.6cm]{\includesvg{figs/03}}

\fourchoices[answer=C]
{合上开关，a 先亮，b 后亮；断开开关，a、b 同时熄灭}
{合上开关，b 先亮，a 后亮；断开开关，a 先熄灭，b 后熄灭}
{合上开关，b 先亮，a 后亮；断开开关，a、b 同时熄灭}
{合上开关，a、b 同时亮；断开开关，b 先熄灭，a 后熄灭}

%% answer: C
%% memo:
\memoanswer{
由于自感作用，合上开关时，b 立即亮，而 a 所在支路开始时电流受到的阻碍很大，电流值很小，而后电流正常，故 a 后亮；断开开关，a、b 和 $L$ 组成一个闭合回路，电路中电流处处相等，稍后同时熄灭，故 C 正确。
}

\item
%% number: 4
%% paperName: 2007年海南理综第 4 题 3 分
%% typeId: 1
%% type: 单选题
%% chapter: 磁场
%% point: 带电粒子在磁场中的运动
%% method: 无
%% score: 3
%% degree: 650
%% duplicateId: 0
%% body:
粒子甲的质量与电荷量分别是粒子乙的 4 倍与 2 倍，两粒子均带正电，让它们在匀强磁场中同一点以大小相等、方向相反的速度开始运动。已知磁场方向垂直纸面向里。以下四个图中，能正确表示两粒子运动轨迹的是\xzanswer{A}

\fourchoices[answer=A, ispicture=true, h=2.4cm]
{figs/04a.png}
{figs/04b.png}
{figs/04c.png}
{figs/04d.png}

%% answer: A
%% memo:
\memoanswer{
带电粒子在磁场中受到的洛伦兹力等于做圆周运动的向心力，有 $qvB=m\frac{v^{2}}{R}$，得 $R=\frac{mv}{qB}$，在速度大小相同的情况下，甲的半径大，故 C、D 错误；洛伦兹力的方向指向圆心，结合左手定则，故 A 正确。
}

\item
%% number: 5
%% paperName: 2007年海南理综第 5 题 3 分
%% typeId: 1
%% type: 单选题
%% chapter: 电路
%% point: 电功电功率
%% method: 无
%% score: 3
%% degree: 650
%% duplicateId: 0
%% body:
一白炽灯泡的额定功率与额定电压分别为 $36\UW$ 与 $36\UV$。若把此灯泡接到输出电压为 $18\UV$ 的电源两端，则灯泡消耗的电功率\xzanswer{B}

\fourchoices[answer=B]
{等于 $36\UW$}
{小于 $36\UW$，大于 $9\UW$}
{等于 $9\UW$}
{小于 $9\UW$}

%% answer: B
%% memo:
\memoanswer{
在额定电压下，灯泡的电阻 $R=\frac{U^{2}}{P}=36\UO$，如果电压降低，功率变小，则灯丝温度下降，电阻变小，加 $18\UV$ 电压时，灯泡的功率 $P=\frac{U^{2}}{R}>9\UW$，故 B 正确。
}

\item
%% number: 6
%% paperName: 2007年海南理综第 6 题 3 分
%% typeId: 1
%% type: 单选题
%% chapter: 电场
%% point: 带电粒子的偏转
%% method: 无
%% score: 3
%% degree: 650
%% duplicateId: 0
%% body:
一平行板电容器中存在匀强电场，电场沿竖直方向。两个比荷（即粒子的电荷量与质量之比）不同的带正电的粒子 a 和 b，从电容器边缘的 P 点（如图）以相同的水平速度射入两平行板之间。测得 a 和 b 与电容器的撞击点到入射点之间的水平距离之比为 $1:2$。若不计重力，则 a 和 b 的比荷之比是\xzanswer{D}

\onepicture[w=4.2cm]{\includesvg{figs/06}}

\fourchoices[answer=D]
{$1:2$}
{$1:1$}
{$2:1$}
{$4:1$}

%% answer: D
%% memo:
\memoanswer{
带电粒子受到的电场力 $F=qE$，产生的加速度为 $a=\frac{F}{m}=\frac{qE}{m}$，在电场中做类平抛运动的时间有 $t=\sqrt{\frac{2d}{a}}$，位移 $x=v_{0}t$，故 $\frac{x_{1}}{x_{2}}=\frac{\sqrt{\frac{m_{1}}{q_{1}}}}{\sqrt{\frac{m_{2}}{q_{2}}}}$，则 $\frac{\frac{q_{1}}{m_{1}}}{\frac{q_{2}}{m_{2}}}=\frac{x_{2}^{2}}{x_{1}^{2}}=\frac{4}{1}$，故 D 正确。
}

\item
%% number: 7
%% paperName: 2007年海南理综第 7 题 4 分
%% typeId: 2
%% type: 多选题
%% chapter: 电场
%% point: 电势能电势差电场力做功
%% method: 无
%% score: 4
%% degree: 650
%% duplicateId: 0
%% body:
如图所示，固定在 Q 点的正点电荷的电场中有 M、N 两点，已知 MQ $<$ NQ。下列叙述正确的是\xzanswer{AD}

\onepicture[w=3.4cm]{\includesvg{figs/07}}

\fourchoices[answer=AD]
{若把一正的点电荷从 M 点沿直线移到 N 点，则电场力对该电荷做功，电势能减少}
{若把一正的点电荷从 M 点沿直线移到 N 点，则该电荷克服电场力做功，电势能增加}
{若把一负的点电荷从 M 点沿直线移到 N 点，则电场力对该电荷做功，电势能减少}
{若把一负的点电荷从 M 点沿直线移到 N 点，再从 N 点沿不同路径移回到 M 点，则该电荷克服电场力做的功等于电场力对该电荷所做的功，电势能不变}

%% answer: AD
%% memo:
\memoanswer{
N 点的电势低于 M 点的电势，正电荷从高电势向低电势运动，电场力做正功，电势能减少，故 A 正确、B 不正确；负电荷从高电势向低电势运动，电场力做负功，电势能增加，故 C 错误；电场力做功与路径无关，只与初末位置有关，故 D 正确。
}

\item
%% number: 8
%% paperName: 2007年海南理综第 8 题 4 分
%% typeId: 2
%% type: 多选题
%% chapter: 直线运动
%% point: 运动图象
%% method: 无
%% score: 4
%% degree: 650
%% duplicateId: 0
%% body:
两辆游戏赛车 a、b 在两条平行的直车道上行驶。$t=0$ 时两车都在同一计时处，此时比赛开始。它们在四次比赛中的 $v-t$ 图如图所示。哪些图对应的比赛中，有一辆赛车追上了另一辆？\xzanswer{AC}

\fourchoices[answer=AC, ispicture=true, h=2.4cm]
{figs/08a.png}
{figs/08b.png}
{figs/08c.png}
{figs/08d.png}

%% answer: AC
%% memo:
\memoanswer{
在 $v-t$ 图象中，图线与 $t$ 轴所夹的面积表示物体发生的位移，A、C 中图线与轴所夹的面积可能相等，B、D 不可能相等，故 A、C 正确，B、D 错误。
}

\item
%% number: 9
%% paperName: 2007年海南理综第 9 题 4 分
%% typeId: 2
%% type: 多选题
%% chapter: 功和能
%% point: 功能原理
%% method: 无
%% score: 4
%% degree: 650
%% duplicateId: 0
%% body:
如图，卷扬机的绳索通过定滑轮用力 $F$ 拉位于粗糙斜面上的木箱，使之沿斜面加速向上移动。在移动过程中，下列说法正确的是\xzanswer{CD}

\onepicture[w=6.5cm]{\includesvg{figs/09}}

\fourchoices[answer=CD]
{$F$ 对木箱做的功等于木箱增加的动能与木箱克服摩擦力所做的功之和}
{$F$ 对木箱做的功等于木箱克服摩擦力和克服重力所做的功之和}
{木箱克服重力做的功等于木箱增加的重力势能}
{$F$ 对木箱做的功等于木箱增加的机械能与木箱克服摩擦力做的功之和}

%% answer: CD
%% memo:
\memoanswer{
由动能定理得 $-W_{G}-W_{f}+W_{F}=\frac{1}{2}mv^{2}-0$，$W_{F}=\frac{1}{2}mv^{2}+W_{G}+W_{f}$，故 A、B 错误；克服重力做的功等于物体重力势能的增加，$\Delta E_{p}=W_{G}$，故 C 正确；又有 $W_{F}=\frac{1}{2}mv^{2}+W_{G}+W_{f}=(\Delta E_{k}+\Delta E_{p})+W_{f}$，故 D 正确。
}

\item
%% number: 10
%% paperName: 2007年海南理综第 10 题 4 分
%% typeId: 2
%% type: 多选题
%% chapter: 运动和力
%% point: 超重失重
%% method: 无
%% score: 4
%% degree: 650
%% duplicateId: 0
%% body:
游乐园中，乘客乘坐能加速或减速运动的升降机，可以体会超重或失重的感觉，下列描述正确的是\xzanswer{BC}

\fourchoices[answer=BC]
{当升降机加速上升时，游客是处在失重状态}
{当升降机减速下降时，游客是处在超重状态}
{当升降机减速上升时，游客是处在失重状态}
{当升降机加速下降时，游客是处在超重状态}

%% answer: BC
%% memo:
\memoanswer{
当升降机加速上升或减速下降时，加速度向上，物体受到的合力向上，支持力大于重力，游客处于超重状态，故 B 正确、A 错误；当升降机加速下降或减速上升时，加速度向下，物体受到的合力向下，支持力小于重力，游客处于失重状态，故 C 正确、D 错误。
}

\item
%% number: 11
%% paperName: 2007年海南理综第 11 题 4 分
%% typeId: 3
%% type: 填空题
%% chapter: 曲线运动
%% point: 天体运动
%% method: 无
%% score: 4
%% degree: 650
%% duplicateId: 0
%% body:
设地球绕太阳做匀速圆周运动，半径为 $R$，速率为 $v$，则太阳的质量可用 $v$、$R$ 和引力常量 $G$ 表示为\tkanswer[1.8cm]{$\frac{v^{2}R}{G}$}。太阳围绕银河系中心的运动可视为匀速圆周运动，其运动速率约为地球公转速率的 7 倍，轨道半径约为地球公转轨道半径的 $2\times10^{9}$ 倍。为了粗略估算银河系中恒星的数目，可认为银河系中所有恒星的质量都集中在银河系中心，且银河系中恒星的平均质量约等于太阳质量，则银河系中恒星数目约为\tkanswer[1.8cm]{$10^{11}$}。

%% answer:
\jdanswer{
$\frac{v^{2}R}{G}$；$10^{11}$
}

%% memo:
\memoanswer{
地球绕太阳做匀速圆周运动，万有引力提供向心力有 $G\frac{Mm}{R^{2}}=m\frac{v^{2}}{R}$，则 $M=\frac{v^{2}R}{G}$；设银河系中的恒星有 $N$ 个，$G\frac{NMM}{(2\times10^{9}R)^{2}}=M\frac{(7v)^{2}}{2\times10^{9}R}$，则 $N=10^{11}$。
}

\item
%% number: 12
%% paperName: 2007年海南理综第 12 题 4 分
%% typeId: 3
%% type: 填空题
%% chapter: 交流电
%% point: 变压器
%% method: 无
%% score: 4
%% degree: 650
%% duplicateId: 0
%% body:
某发电厂用 $2.2\UkV$ 的电压将电能输送到远处的用户，后改用 $22\UkV$ 的电压，在既有输电线路上输送同样的电功率。前后两种输电方式消耗在输电线上电功率之比为\tkanswer[1.5cm]{100}。要将 $2.2\UkV$ 的电压升高至 $22\UkV$，若变压器原线圈的匝数为 180 匝，则副线圈的匝数应该是\tkanswer[1.5cm]{1800}匝。

%% answer:
\jdanswer{
$100$；$1800$
}

%% memo:
\memoanswer{
由于输送功率相同，$\frac{I_{1}}{I_{2}}=\frac{U_{2}}{U_{1}}=\frac{10}{1}$，电路中损耗的功率 $P=I^{2}R$，损耗的功率 $\frac{P_{1}}{P_{2}}=\frac{100}{1}$；匝数之比等于电压之比，$\frac{n_{1}}{n_{2}}=\frac{U_{1}}{U_{2}}=\frac{180}{1800}$。
}

\item
%% number: 13
%% paperName: 2007年海南理综第 13 题 7 分
%% typeId: 4
%% type: 实验题
%% chapter: 电路
%% point: 电路实验
%% method: 无
%% score: 7
%% degree: 650
%% duplicateId: 0
%% body:
图（a）中电源电动势为 $E$，内阻可忽略不计；电流表 mA 具有一定的内阻，电压表 V 的内阻不是无限大，S 为单刀双掷开关，$R$ 为待测电阻。当 S 向电压表一侧闭合时，电压表读数为 $U_{1}$，电流表读数为 $I_{1}$；当 S 向 $R$ 一侧闭合时，电流表读数为 $I_{2}$。

\begin{enumerate}
	\item
	根据已知条件与测量数据，可以得出待测电阻 $R=$\tkanswer[2.5cm]{$\frac{E}{I_{2}}-\frac{E-U_{1}}{I_{1}}$}。
	\item
	根据图（a）所给出的电路，在图（b）的各器件实物图之间画出连接的导线。
\end{enumerate}

\twopicture[align=right, labelprefix={}, numA={（a）}, numB={（b）}, wA=5.2cm, wB=5.0cm]
{figs/13a.png}
{figs/13b.png}

%% answer:
\jdanswer{
\begin{enumerate}
	\item
	$R=\frac{E}{I_{2}}-\frac{E-U_{1}}{I_{1}}$。
	\item
	连线如图。

	\onepicture[w=4.6cm]{figs/13_an.png}
\end{enumerate}
}

%% memo:
\memoanswer{
（1）当 S 向电压表一侧闭合时，电压表读数为 $U_{1}$，电流表读数为 $I_{1}$，设电流表的内阻为 $r$，由闭合电路的欧姆定律得 $E=U_{1}+I_{1}r$，整理得 $r=\frac{E-U_{1}}{I_{1}}$；当 S 向 $R$ 一侧闭合时，电流表读数为 $I_{2}$，由闭合电路的欧姆定律得 $E=I_{2}(R+r)$，整理得 $R=\frac{E}{I_{2}}-r=\frac{E}{I_{2}}-\frac{E-U_{1}}{I_{1}}$。

（2）按照电路图连线，要注意电流从电压表和电流表的正接线柱流入、负接线柱流出，单刀双掷开关控制整个电路的通断，连线不要交叉。
}

\item
%% number: 14
%% paperName: 2007年海南理综第 14 题 8 分
%% typeId: 4
%% type: 实验题
%% chapter: 运动和力
%% point: 牛顿第二定律
%% method: 无
%% score: 8
%% degree: 450
%% duplicateId: 0
%% body:
现要验证“当质量一定时，物体运动的加速度与它所受的合外力成正比”这一物理规律。给定的器材如下：一倾角可以调节的长斜面（如图）、小车、计时器一个、米尺。

\begin{enumerate}
	\item
	填入适当的公式或文字，完善以下实验步骤（不考虑摩擦力的影响）：

	\begin{steps}
		\item
		让小车自斜面上方一固定点 $A_{1}$ 从静止开始下滑到斜面底端 $A_{2}$，记下所用的时间 $t$。
		\item
		用米尺测量 $A_{1}$ 与 $A_{2}$ 之间的距离 $s$，则小车的加速度 $a=$\tkanswer[2cm]{$\frac{2s}{t^{2}}$}。
		\item
		用米尺测量 $A_{1}$ 相对于 $A_{2}$ 的高度 $h$。设小车所受重力为 $mg$，则小车所受的合外力 $F=$\tkanswer[2cm]{$\frac{mgh}{s}$}。
		\item
		改变\tkanswer[2.6cm]{斜面倾角（或填 $h$ 的数值）}，重复上述测量。
		\item
		以 $h$ 为横坐标，$1/t^{2}$ 为纵坐标，根据实验数据作图。如能得到一条过原点的直线，则可验证“当质量一定时，物体运动的加速度与它所受的合外力成正比”这一规律。
	\end{steps}

	\item
	在探究如何消除上述实验中摩擦阻力影响的过程中，某同学设计的方案是：

	\begin{steps}
		\item
		调节斜面倾角，使小车在斜面上匀速下滑。测量此时 $A_{1}$ 点相对于斜面底端 $A_{2}$ 的高度 $h_{0}$。
		\item
		进行（1）中的各项测量。
		\item
		计算与作图时用 $(h-h_{0})$ 代替 $h$。
	\end{steps}

	对此方案有以下几种评论意见：

	（A）方案正确可行。

	（B）方案的理论依据正确，但利用所给的器材无法确定小车在斜面上是否做匀速运动。

	（C）方案的理论依据有问题，小车所受摩擦力与斜面倾角有关。

	其中合理的意见是\tkanswer[1.5cm]{C}。
\end{enumerate}

\onepicture[w=5.8cm, align=right]{\includesvg{figs/14}}

%% answer:
\jdanswer{
\begin{enumerate}
	\item
	$\frac{2s}{t^{2}}$；$\frac{mgh}{s}$；斜面倾角（或填 $h$ 的数值）。
	\item
	C。
\end{enumerate}
}

%% memo:
\memoanswer{
本题原卷及材料中未提供详解，待补充。
}

\item
%% number: 15
%% paperName: 2007年海南理综第 15 题 8 分
%% typeId: 6
%% type: 计算题
%% chapter: 磁场
%% point: 安培力
%% method: 无
%% score: 8
%% degree: 650
%% duplicateId: 0
%% body:
据报道，最近已研制出一种可以投入使用的电磁轨道炮，其原理如图所示。炮弹（可视为长方形导体）置于两固定的平行导轨之间，并与轨道壁密接。开始时炮弹在轨道的一端，通以电流后炮弹会被磁力加速，最后从位于导轨另一端的出口高速射出。设两导轨之间的距离 $w=0.10\Um$，导轨长 $L=5.0\Um$，炮弹质量 $m=0.30\Ukg$。导轨上的电流 $I$ 的方向如图中箭头所示。可认为，炮弹在轨道内运动时，它所在处磁场的磁感应强度始终为 $B=2.0\UT$，方向垂直于纸面向里。若炮弹出口速度为 $v=2.0\times10^{3}\Ums$，求通过导轨的电流 $I$。忽略摩擦力与重力的影响。

\onepicture[w=5.5cm, align=right]{\includesvg{figs/15}}

%% answer:
\jdanswer{
$I=6.0\times10^{5}\UA$
}

%% memo:
\memoanswer{
在导轨通有电流 $I$ 时，炮弹作为导体受到磁场施加的安培力为 $F=BIw$①，设炮弹的加速度的大小为 $a$，则有 $F=ma$②，炮弹在两导轨间做匀加速运动，因而 $v^{2}=2aL$③，联立①②③式得 $I=\frac{\frac{1}{2}mv^{2}}{BwL}$④，代入题给数据得 $I=6.0\times10^{5}\UA$⑤。
}

\item
%% number: 16
%% paperName: 2007年海南理综第 16 题 11 分
%% typeId: 6
%% type: 计算题
%% chapter: 运动和力
%% point: 牛顿定律的应用
%% method: 整体与隔离
%% score: 11
%% degree: 650
%% duplicateId: 0
%% body:
如图所示，一辆汽车 A 拉着装有集装箱的拖车 B，以速度 $v_{1}=30\Ums$ 进入向下倾斜的直车道。车道每 $100\Um$ 下降 $2\Um$。为了使汽车速度在 $s=200\Um$ 的距离内减到 $v_{2}=10\Ums$，驾驶员必须刹车。假定刹车时地面的摩擦阻力是恒力，且该力的 70％ 作用于拖车 B，30％ 作用于汽车 A。已知 A 的质量 $m_{1}=2000\Ukg$，B 的质量 $m_{2}=6000\Ukg$。求汽车与拖车的连接处沿运动方向的相互作用力。取重力加速度 $g=10\Umsq$。

\onepicture[w=6.5cm, align=right]{\includesvg{figs/16}}

%% answer:
\jdanswer{
$f_{N}=880\UN$
}

%% memo:
\memoanswer{
汽车沿倾斜车道做匀减速运动，用 $a$ 表示加速度的大小，有 $v_{2}^{2}-v_{1}^{2}=-2as$①，用 $F$ 表示刹车时的阻力，由牛顿第二定律得 $F-(m_{1}+m_{2})g\sin\alpha=(m_{1}+m_{2})a$②，式中 $\sin\alpha=\frac{2}{100}=2\times10^{-2}$③，设刹车过程中地面作用于汽车的阻力为 $f$，由题意得 $f=\frac{30}{100}F$④，方向与汽车前进方向相反，用 $f_{N}$ 表示拖车作用于汽车的力，设其方向与汽车前进方向相同。以汽车为研究对象，由牛顿第二定律得 $f-f_{N}-m_{1}g\sin\alpha=m_{1}a$⑤，由②④⑤式得 $f_{N}=\frac{30}{100}(m_{1}+m_{2})(a+g\sin\alpha)-m_{1}(a+g\sin\alpha)$⑥，由①③⑥式，代入有关数据得 $f_{N}=880\UN$。
}

\item
%% number: 17
%% paperName: 2007年海南理综第 17 题 8 分
%% typeId: 6
%% type: 计算题
%% chapter: 气体性质
%% point: 理想气体状态方程
%% method: 临界
%% score: 8
%% degree: 650
%% duplicateId: 0
%% body:
\begin{enumerate}
	\item
	（4 分）有以下说法：

	\onepicture[w=3.4cm]{\includesvg{figs/17a}}

	\eightchoices[answer=ABEG]
	{气体的温度越高，分子的平均动能越大}
	{即使气体的温度很高，仍有一些分子的运动速率是非常小的}
	{对物体做功不可能使物体的温度升高}
	{如果气体分子间的相互作用力小到可以忽略不计，则气体的内能只与温度有关}
	{一由不导热的器壁做成的容器，被不导热的隔板分成甲、乙两室。甲室中装有一定质量的温度为 $T$ 的气体，乙室为真空，如图所示。提起隔板，让甲室中的气体进入乙室。若甲室中的气体的内能只与温度有关，则提起隔板后当气体重新达到平衡时，其温度仍为 $T$}
	{空调机作为制冷机使用时，将热量从温度较低的室内送到温度较高的室外，所以制冷机的工作是不遵守热力学第二定律的}
	{对于一定量的气体，当其温度降低时，速率大的分子数目减少，速率小的分子数目增加}
	{从单一热源吸取热量使之全部变成有用的机械功是不可能的}

	其中正确的是\tkanswer[1.5cm]{ABEG}。（选错一个扣 1 分，选错两个扣 3 分，选错三个或三个以上得 0 分，最低得分为 0 分）

	\item
	（8 分）如图，在大气中有一水平放置的固定圆筒，它由 a、b 和 c 三个粗细不同的部分连接而成，各部分的横截面积分别为 $2S$、$\frac{1}{2}S$ 和 $S$。已知大气压强为 $p_{0}$，温度为 $T_{0}$。两活塞 A 和 B 用一根长为 $4l$ 的不可伸长的轻线相连，把温度为 $T_{0}$ 的空气密封在两活塞之间，此时两活塞的位置如图所示。现对被密封的气体加热，使其温度缓慢上升到 $T$。若活塞与圆筒壁之间的摩擦可忽略，此时两活塞之间气体的压强可能为多少？

	\onepicture[w=6.0cm, align=right]{\includesvg{figs/17b}}
\end{enumerate}

%% answer:
\jdanswer{
\begin{enumerate}
	\item
	$ABEG$。
	\item
	当 $T\leqslant\frac{5}{4}T_{0}$ 时，气体的压强为 $p_{2}=p_{0}$；当 $T>\frac{5}{4}T_{0}$ 时，气体的压强为 $p=\frac{4T}{5T_{0}}p_{0}$。
\end{enumerate}
}

%% memo:
\memoanswer{
（1）温度是分子平均动能的标志，气体的温度越高，分子的平均动能越大，故 A 正确；由于分子的运动是无规则的、杂乱无章的，所以即使气体的温度很高，仍有一些分子的运动速率是非常小的，故 B 正确；做功是改变物体内能的一种方式，做功可以使物体的温度升高，故 C 错误；如果气体分子间的相互作用力忽略不计，分子势能不计，只有分子动能，气体的内能只与分子动能有关，而分子平均动能与温度有关，故气体的内能应与分子数和温度有关，故 D 错误；因为在去除隔板时，气体没有对外做功，容器又不导热，所以由热力学第一定律可知，气体的内能不变，由题，甲室中气体的内能只与温度有关，故甲室中气体的温度不变，故 E 正确；空调机作为制冷机使用时，在压缩机做功的情况下，将热量从温度较低的室内送到温度较高的室外，并不违反热力学第二定律，故 F 错误；当气体温度降低时，分子的平均速率减小，说明速度大的分子数目减少，速率小的分子数目增加，故 G 正确；由热力学第二定律可知，在没有外界影响之下，从单一热源吸取热量使之全部变成有用的机械功是不可能的，故 H 错误。

（2）设加热前，被密封气体的压强为 $p_{1}$，轻线的张力为 $f$。因两活塞都处在静止状态，对活塞 A，有 $2p_{0}S-2p_{1}S+f=0$①，对活塞 B，有 $p_{1}S-p_{0}S-f=0$②，由①②式得 $p_{1}=p_{0}$③，$f=0$④。即被密封气体的压强与大气压强相等，细线处在拉直的松弛状态。这时气体的体积为 $V_{1}=2Sl+Sl+Sl=4Sl$⑤，对气体加热时，被密封气体温度缓慢升高，两活塞一起向左缓慢移动，气体体积增大，压强保持 $p_{1}$ 不变，若持续加热，此过程会一直持续到活塞向左移动的距离等于 $l$ 为止，这时气体的体积为 $V_{2}=4Sl+Sl=5Sl$⑥，设此时气体的温度为 $T_{2}$，由盖-吕萨克定律得 $\frac{V_{2}}{T_{2}}=\frac{V_{1}}{T_{0}}$⑦，由③⑥⑦式得 $T_{2}=\frac{5}{4}T_{0}$⑧，由此可知，当 $T\leqslant T_{2}=\frac{5}{4}T_{0}$ 时，气体的压强 $p_{2}=p_{0}$⑨，当 $T>T_{2}$ 时，活塞已无法移动，被密封气体的体积保持 $V_{2}$ 不变，气体经历一等容升压过程，当气体的温度为 $T$ 时，设其压强为 $p$，由查理定律得 $\frac{p}{T}=\frac{p_{0}}{T_{2}}$⑩，由⑧⑨⑩式得 $p=\frac{4T}{5T_{0}}p_{0}$，即当 $T>\frac{5}{4}T_{0}$ 时，气体的压强为 $\frac{4T}{5T_{0}}p_{0}$。
}

\item
%% number: 18
%% paperName: 2007年海南理综第 18 题 7 分
%% typeId: 6
%% type: 计算题
%% chapter: 光学
%% point: 光的折射
%% method: 无
%% score: 7
%% degree: 450
%% duplicateId: 0
%% body:
\begin{enumerate}
	\item
	（5 分）一列简谐横波，沿 $x$ 轴正向传播，位于原点的质点的振动图象如图（a）所示。

	\twopicture[labelprefix={}, numA={（a）}, numB={（b）}, wA=6.0cm, wB=6.0cm]
	{figs/18a.png}
	{figs/18b.png}

	①该振动的振幅是\tkanswer[1.2cm]{8}$\Ucm$；

	②振动的周期是\tkanswer[1.2cm]{0.2}$\Us$；

	③在 $t$ 等于 $1/4$ 周期时，位于原点的质点离开平衡位置的位移是\tkanswer[1.2cm]{0}$\Ucm$。

	图（b）为该波在某一时刻的波形图，A 点位于 $x=0.5\Um$ 处。

	④该波的传播速度是\tkanswer[1.5cm]{10}$\Ums$；

	⑤经过 $1/2$ 周期后，A 点离开平衡位置的位移是\tkanswer[1.5cm]{$-8$}$\Ucm$。

	\item
	（7 分）如图，置于空气中的一不透明容器内盛满某种透明液体。容器底部靠近器壁处有一竖直放置的 $6.0\Ucm$ 长的线光源。靠近线光源一侧的液面上盖有一遮光板，另一侧有一水平放置的与液面等高的望远镜，用来观察线光源。开始时通过望远镜不能看到线光源的任何一部分。将线光源沿容器底向望远镜一侧平移至某处时，通过望远镜刚好可以看到线光源底端，再将线光源沿同一方向移动 $8.0\Ucm$，刚好可以看到其顶端。求此液体的折射率 $n$。

	\onepicture[w=5.8cm, align=right]{\includesvg{figs/18c}}
\end{enumerate}

%% answer:
\jdanswer{
\begin{enumerate}
	\item
	① $8$；② $0.2$；③ $0$；④ $10$；⑤ $-8$。
	\item
	$n=1.3$。
\end{enumerate}
}

%% memo:
\memoanswer{
（1）由振动图象可以看出该振动的振幅为 $A=8\Ucm$，振动周期为 $T=0.2\Us$；在 $t$ 等于 $\frac{1}{4}$ 周期时，位于原点的质点刚好回到平衡位置，因而位移为 0；由图乙可以看出，该波的波长为 $\lambda=2\Um$，则波速 $v=\frac{\lambda}{T}=\frac{2}{0.2}\Ums=10\Ums$；经过 $\frac{1}{2}$ 周期后，A 点刚好到达负的最大位移处，因而位移为 $-8\Ucm$。

（2）当线光源上某一点发出的光线射到未被遮光板遮住的液面上时，射到遮光板边缘 O 的那条光线的入射角最小。若线光源底端在 A 点时，望远镜内刚好可以看到此光源底端，设过 O 点液面的法线为 $OO_{1}$，则 $\angle AOO_{1}=\alpha$①，其中 $\alpha$ 为此液体到空气的全反射临界角。由折射定律有 $\sin\alpha=\frac{1}{n}$②，同理，若线光源顶端在 $B_{1}$ 点时，通过望远镜刚好可以看到此光源顶端，则 $\angle B_{1}OO_{1}=\alpha$，设此时线光源底端位于 B 点，由图中几何关系可得 $\sin\alpha=\frac{\overline{AB}}{\overline{AB_{1}}}$③，联立②③式得 $n=\frac{\sqrt{\overline{AB}^{2}+\overline{BB_{1}}^{2}}}{\overline{AB}}$④，由题给条件可知 $\overline{AB}=8.0\Ucm$，$\overline{BB_{1}}=6.0\Ucm$，代入④式得 $n=1.3$⑤。

\onepicture[w=5.5cm]{\includesvg{figs/18_an}}
}

\item
%% number: 19
%% paperName: 2007年海南理综第 19 题 6 分
%% typeId: 6
%% type: 计算题
%% chapter: 运动和力
%% point: 动量守恒定律
%% method: 无
%% score: 6
%% degree: 450
%% duplicateId: 0
%% body:
\begin{enumerate}
	\item
	（6 分）氢原子第 $n$ 能级的能量为 $E_{n}=E_{1}/n^{2}$，其中 $E_{1}$ 是基态能量。而 $n=1$，2，$\ldots$。若一氢原子发射能量为 $-\frac{3}{16}E_{1}$ 的光子后处于比基态能量高出 $-\frac{3}{4}E_{1}$ 的激发态，则氢原子发射光子前后分别处于第几能级？
	\item
	（6 分）一速度为 $v$ 的高速 $\alpha$ 粒子（$\ce{^{4}_{2}He}$）与同方向运动的氖核（$\ce{^{20}_{10}Ne}$）发生弹性正碰，碰后 $\alpha$ 粒子恰好静止。求碰撞前后氖核的速度（不计相对论修正）
\end{enumerate}

%% answer:
\jdanswer{
\begin{enumerate}
	\item
	第 4 与第 2 能级。
	\item
	$v_{0}=\frac{2}{5}v$；$v_{\mathrm{Ne}}=\frac{3}{5}v$。
\end{enumerate}
}

%% memo:
\memoanswer{
（1）设氢原子发射光子前后分别处于第 $l$ 与第 $m$ 能级，则由题意得 $\frac{E_{1}}{l^{2}}-\frac{E_{1}}{m^{2}}=-\frac{3}{16}E_{1}$①，$\frac{E_{1}}{m^{2}}-E_{1}=-\frac{3}{4}E_{1}$②，由②式解得 $m=2$③，由①③式得 $l=4$④，氢原子发射光子前后分别处于第 4 与第 2 能级。

（2）设 $\alpha$ 粒子与氖核的质量分别为 $m_{\alpha}$ 与 $m_{\mathrm{Ne}}$，氖核在碰撞前后的速度分别为 $v_{0}$ 与 $v_{\mathrm{Ne}}$。由动量守恒与机械能守恒得 $m_{\alpha}v+m_{\mathrm{Ne}}v_{0}=m_{\mathrm{Ne}}v_{\mathrm{Ne}}$①，$\frac{1}{2}m_{\alpha}v^{2}+\frac{1}{2}m_{\mathrm{Ne}}v_{0}^{2}=\frac{1}{2}m_{\mathrm{Ne}}v_{\mathrm{Ne}}^{2}$②，解得 $v_{0}=\frac{m_{\mathrm{Ne}}-m_{\alpha}}{2m_{\mathrm{Ne}}}v$③，$v_{\mathrm{Ne}}=\frac{m_{\mathrm{Ne}}+m_{\alpha}}{m_{\mathrm{Ne}}}v$④，已知 $\frac{m_{\alpha}}{m_{\mathrm{Ne}}}=\frac{1}{5}$⑤，将⑤式代入③④式得 $v_{0}=\frac{2}{5}v$⑥，$v_{\mathrm{Ne}}=\frac{3}{5}v$⑦。
}

\end{enumerate}

\end{document}
